\documentclass[10pt]{article}
\usepackage{titling}

\title{Nucleophilic Substitution and Elimination}
\date{DATE PLACEHOLDER}
\author{Ingyu Bahng}

\usepackage[letterpaper, margin=1in]{geometry}
\usepackage{amsmath} % better math functions
\usepackage{amssymb} % better math symbols
\usepackage{babel}[english] % change rules to english
\usepackage{lipsum} % allows for lorem ipsum text generation
\usepackage{xr-hyper} % allows cross referencing labels from external LaTeX documents 
\usepackage{hyperref} % for hyperlinking text
\usepackage{fancyvrb} % for better verbatim environments
\usepackage{newverbs} % allows for inline typewriter font via \texttt{}
\usepackage{fancyvrb} % also code font blocks but has more capabilities than texttt
\usepackage{footnote} % allows footnotes inside tabulars
\usepackage{dcolumn} % allows for decimal alignment in tables
\usepackage{tabularx}

\usepackage{fontspec}
\setmainfont{Gentium Plus} % allowing greek letters in text

\usepackage{unicode-math}
\setmathfont{XITS Math} % allowing greek letters in math



% tcolorbox package
\usepackage{tcolorbox}
\tcbuselibrary{breakable}

\tcbset{ % this is the default design of tcolorboxes
  colframe = black,
  colback  = black!1,
  coltitle = black,
  colbacktitle = black!15,
  breakable, 
  arc=0mm,
  boxrule=0.5pt,
  toptitle=3pt,
  bottomtitle=3pt,
  left=8pt,
  right=8pt,
  top=6pt,
  bottom=6pt,
  before skip=12pt,
  after skip=12pt,
  bottomrule at break=-1pt,
  toprule at break=-1pt,
  fonttitle=\bfseries,
}

% Custom Environments
\newtcolorbox[auto counter, number within=section]{definition}[1][]
{
    title = \textbf{Definition \thetcbcounter: ~#1}
}

\newtcolorbox[auto counter, number within=section]{core}[1][]
{
    title = \textbf{Core Concept \thetcbcounter: ~#1}
}


\usepackage{fancyhdr}
\usepackage[skip=0.8em, indent=0em]{parskip}

\pagestyle{fancy} % this section generates the lines and text at the top and bottom of each page
\fancyhead[L]{\thetitle}
\fancyhead[C]{\theauthor}
\fancyhead[R]{\thedate}
\fancyfoot[C]{\thepage}
\renewcommand{\footrulewidth}{0.3pt}
\renewcommand{\thispagestyle}[1]{}

\renewcommand{\arraystretch}{1.5} % table vertical padding
\setlength{\tabcolsep}{10pt} % table horizontal padding

\usepackage{enumitem} % better control over enumerate and itemize

% List customization
\setlist[itemize]{%
  label=\bullet,
  itemsep=2pt,
  leftmargin=1.5em,
}

\setlist[enumerate]{%
  label=(\alph*),
  itemsep=2pt,
  leftmargin=1.5em,
}



\usepackage{pgfplots} % for plot creation
\pgfplotsset{compat=1.18} % setting the version used for pgfplots
\usepackage{float} % for better figure placement using [h]
\usepackage{graphicx} % for importing figures into the tex file
\usepackage{subcaption} % allows for multiple subfigures with individual captions wihtin one figure
\usepackage{xparse} % allows for more than 9 parameters in commands
\usepackage{adjustbox} % for valign option
\captionsetup{font=small} % setting the caption fontsize to small always

\usepackage{silence}
\WarningsOff[float]  % suppress float package warnings

\newcommand{\myfig}[4]{%
  \begin{figure}[H]
    \centering
    \adjustbox{valign=c}{\includegraphics[width=#1\textwidth]{#2}}
    \caption{#3}
    \label{#4}
  \end{figure}
}

\newcommand{\mymultifig}[8]{
  \begin{figure}[H]%
    \centering%
    \begin{subfigure}{#1\textwidth}%
      \centering%
      \includegraphics[width=\linewidth]{#2}%
    \end{subfigure}%
    \hfill%
    \begin{subfigure}{#3\textwidth}%
      \centering%
      \includegraphics[width=\linewidth]{#4}%
    \end{subfigure}%
    \hfill%
    \begin{subfigure}{#5\textwidth}%
      \centering%
      \includegraphics[width=\linewidth]{#6}%
    \end{subfigure}%
    \caption{#7}%
    \label{#8}%
  \end{figure}%
}

\usepackage{newverbs} % allows for inline typewriter font via \texttt{}
\usepackage{fancyvrb} % also code font blocks but has more capabilities than texttt
\usepackage{listings} % allows for codeblocks - config below

%BELOW ARE CODE BLOCK DESIGNS (Light Theme)
\definecolor{gray}{HTML}{707070}
\definecolor{lightgray}{HTML}{333333}
\definecolor{codebg}{HTML}{F5F5F5}
\definecolor{keyword}{HTML}{0000FF}
\definecolor{string}{HTML}{A31515}
\definecolor{number}{HTML}{098658}
\definecolor{function}{HTML}{795E26}
\definecolor{comment}{HTML}{008000}
\definecolor{classname}{HTML}{267F99}

% default options for listings (for code)
\lstset{
  frame=ltbr,
  language=python,
  aboveskip=3mm,
  belowskip=3mm,
  showstringspaces=false,
  columns=fullflexible,
  keepspaces=true,
  basicstyle=\fontsize{9}{10}\ttfamily\color{lightgray},
  numbers=left,
  firstnumber=1,
  numberstyle=\fontsize{7}{10}\ttfamily\color{gray},
  keywordstyle=\color{keyword},
  commentstyle=\color{comment},
  stringstyle=\color{string},
  backgroundcolor=\color{codebg},
  breaklines=true,
  breakatwhitespace=true,
  tabsize=2,
  xleftmargin=3em,
  numbersep=1em,
  framexleftmargin=2.5em,
  framesep=6pt,
  stepnumber=1,
}

\hfuzz=5.002pt % ignore overfull hbox badness warnings below this limit



\usepackage{chemfig}
\usepackage[version=4]{mhchem}

\newcommand{\bce}[1]{
  \hspace{7mm}{#1}
}

\newcommand{\sno}[0]{
  $\mathrm{S}_{\mathrm{N}}1$
}

\newcommand{\snt}[0]{
  $\mathrm{S}_{\mathrm{N}}2$
}

\newcommand{\reactionfig}[4]{
    \begin{figure}[H]
    \centering
    \scalebox{#4}{
    \schemestart
    #1
    \schemestop
    }
    \caption{#2}
    \label{#3}
    \end{figure}
}




\externaldocument[gas-]{../gases/gases}[../gases/gases.pdf]
\externaldocument[aci-]{../acids_and_bases/acids_and_bases}[../acids_and_bases/acids_and_bases.pdf]
\externaldocument[equ-]{../equilibrium/equilibrium}[../equilibrium/equilibrium.pdf]
\externaldocument[the-]{../thermodynamics/thermodynamics}[../thermodynamics/thermodynamics.pdf]
\externaldocument[qua-]{../quantization/quantization}[../quantization/quantization.pdf]
\externaldocument[alc-]{../alcohols/alcohols}[../alcohols/alcohols.pdf]
\externaldocument[eth-]{../ethers/ethers}[../ethers/ethers.pdf]
\externaldocument[nom-]{../nomenclature/nomenclature}[../nomenclature/nomenclature.pdf]
\externaldocument[rad-]{../radical_halogenation/radical_halogenation}[../radical_halogenation/radical_halogenation.pdf]
\externaldocument[con-]{../conformational_analysis/conformational_analysis}[../conformational_analysis/conformational_analysis.pdf]
\externaldocument[alk-]{../alkane_alkene_alkyne/alkane_alkene_alkyne}[../alkane_alkene_alkyne/alkane_alkene_alkyne.pdf]
%\externaldocument[nuc-]{../nuc_sub_and_elim/nuc_sub_and_elim}[../nuc_sub_and_elim/nuc_sub_and_elim.pdf]

\begin{document}

\input{../../presets/_general/startup}

\section{Nucleophilic Substitution}

  After generating a haloalkane compound via radical halogenation (\textit{as covered in \href{../radical_halogenation/radical_halogenation.pdf}{Radical Halogenation}}), the next step is to substitute that halogen with a different atom/molecule of our choosing. One of the methods we can use to carry this out is nucleophilic substitution.

  \begin{core}{Nucleophilic Substitution}
   
    \textbf{Nucleophilic substitution} is a fundamental class of organic reactions where an \ce{e^-} rich species, called a \textbf{nucleophile (Nu)}, attacks an electron-deficient atom called the \textbf{electrophile (E)}, displacing a specific group known as the \textbf{leaving group (L)}, which departs with the bonding electron pair. Halogens are common leaving groups due to their electronegativity and ability to stabilise the negative charge upon departure.

  \end{core}

  Nucleophilic substitution can be categorized into \snt and \sno.

  \begin{core}{\snt}
    
    \snt reactions are \textbf{bimolecular}, meaning that both the leaving group and nucleophile are present in the transition state and making the reaction rate dependent on both [Nu] and [E-L] ($\text{Rate} = k[\ce{CH3Cl}][\ce{^-OH}]$). The following is a chemical reaction of \ce{CH3Cl} undergoing \snt reaction with \ce{OH-}.

    \vspace{15pt}
    \reactionfig{
      \scheme{
        \chemfig{
          @{oxygen}\textcolor{Red}{\ce{HO-}}
        }
        \+
        \chemfig{
          [:0]@{carbon}\ce{CH3}
          -[@{bond}:0,,,,Red]@{chlorine}\textcolor{Green}{Cl}
        }
        \arrow{->}
        \chemfig{
          \ce{CH3OH}
        }
        \+
        \chemfig{
          \textcolor{Green}{Cl}
        }
      }
      \chemmovearrow{
        \reactionfigarrow{oxygen}{60}{carbon}{120}{Red}
        \reactionfigarrow{bond}{100}{chlorine}{80}{Red}
      }
    }{\snt equation}{fig:filler1}{1.0}

    Here, Cl is the leaving group and \ce{OH-} is the nucleophile. From a 3D point of view, we see the following outline interaction.

    \reactionfig{
      \scheme{
        \chemname{
          \chemfig{
            @{nucleophile}\textcolor{Red}{\ce{Nu-}}
          }
        }{\textcolor{Red}{\textit{Nucleophile}}}
        \+
        \chemname{
          \chemfig{
            [:0]@{electrophile}\textcolor{Blue}{C^{δ^+}}
            (-[:120])
            (<[:-100])
            (<:[:-160])
            -[@{bond}:0,,,,Red]@{leaving}\textcolor{Green}{X^{δ^-}}
          }
        }{\textcolor{Blue}{\textit{Electrophile}}---\textcolor{Green}{\textit{Leaving Group}}}
        \arrow{->}
        \chemfig{
          [:0]\textcolor{Red}{Nu}
          -[:0,,,,Red]\textcolor{Blue}{C}
          (<[:-80])
          (<:[:-20])
          -[:60]
        }
        \+
        \chemfig{
          \textcolor{Green}{\ce{X-}}
        }
      }
      \chemmovearrow{
        \reactionfigarrow{nucleophile}{60}{electrophile}{180}{Red}
        \reactionfigarrow{bond}{100}{leaving}{80}{Red}
      }
    }{\snt mechanism reaction}{fig:filler2}{0.9}

    It is important to note that \snt reactions occur in \textbf{one step}, where bond making and breaking occur simultaneously, meaning there are no intermediates --- unlike the two step propagation of radical halogenation (EXTERNAL REFERENCE TO RAD HAL). Consequently, the energy diagram consists of a single hill and is exothermic when the nucleophile is a more favorable attachment than the leaving group. \textit{L is denoted as X in the figure below.}

    \pgfig{
      \begin{axis}[
          width=7cm,
          height=7cm,
          xmin=-4, xmax=4,
          ymin=-0.3, ymax=1.4,
          xtick=\empty,
          ytick=\empty,
          xlabel={},
          ylabel={},
          axis line style={thick},
          \axispresets
      ]

      % energy curve — gaussian-like hill
      \addplot[
          thick,
          black,
          domain=-1.8:1.5,
          samples=200,
      ] {exp(-x^2) - 0.15*x};

      % reactant energy level (left)
      \addplot[thick, black] coordinates {(-3.5, 0.30916) (-1.8, 0.30916)};

      % product energy level (right) — slightly lower
      \addplot[thick, black] coordinates {(1.5, -0.1196) (3.5, -0.1196)};

      % transition state dot
      \addplot[mark=*, mark size=2.5pt, black] coordinates {(-0.075, 1.0056)};

      % transition state structure label
      \node[cyan!70!blue, font=\small] at (axis cs:0, 1.15) {%
          $\left[\,\text{Nuc}\cdots\text{C}\cdots\text{X}\,\right]$
      };

      % reactant label
      \node[font=\small] at (axis cs:-2.5, 0.23) {$\text{R}{-}\text{X} + \text{Nuc:}^{-}$};

      % product label
      \node[font=\small] at (axis cs:2.8, -0.2) {$\text{R}{-}\text{Nuc} + \text{X}^{-}$};

      % "single transition state" label
      \node[cyan!70!blue, font=\small, align=center] at (axis cs:2.8, 0.75) {single\\transition\\state};

      \end{axis}
    }{\snt energy diagram}{fig:filler3}{1.0}

    A result of the reaction being one step is \textbf{stereospecificity} --- there is \textit{always} an inversion. A few key consequences of this are worth keeping in mind:

    \begin{enumerate}

      \item Inversion does not always imply a shift between the R and S configuration (EXTERNAL REFERENCE TO CONFORMATIONAL ANALYSIS)

      \item Two successive \snt reactions result in structural retention (double inversion)

      \item Performing \snt on only one chiral center of a molecule with two chiral centers yields a diastereoisomer (a non-mirror-image stereoisomer)

      \item Starting with a \textbf{racemate} (a 50/50 mix of R and S enantiomers) will still yield a racemate, as all configurations are inverted
        
    \end{enumerate}

    All of these points follow intuitively from the structural consequences of the \snt mechanism, though they may not be immediately obvious.

  \end{core}  

  Before exploring \sno reactions, it is worth establishing what defines a \textbf{"good" leaving group and nucleophile}, as this governs the outcome of both mechanisms.

  \begin{core}{Leaving Group Quality}

    There are two factors that determine whether a leaving group is good: \textbf{electronegativity} and \textbf{orbital size}. Regarding electronegativity, a good leaving group is also a weak base \ce{A-} --- one that is stable with an extra \ce{e^-} due to its high electronegativity. Regarding orbital size, a larger orbital size element forms relatively weaker bonds, as its molecular bonding orbitals occur at higher energy levels compared to smaller elements, making it more inclined to dissociate. Together, these two factors mean that better leaving groups appear moving left to right ($\rightarrow$) and top to bottom ($\downarrow$) across the periodic table.

    \pgfig{
      \miniptable

      % PT label
      \node[font=\small] at (2.5, 1.4) {Periodic Table};

      % EN > OS arrow (left to right, along top)
      \draw[->, thick] (0.8, 3.5) -- (4.8, 3.5);
      \node[font=\small] at (2.8, 3.8) {EN $>$ OS};

      % OS > EN arrow (top to bottom, right side)
      \draw[->, thick] (5.5, 3.2) -- (5.5, 0.2);
      \node[font=\small] at (6.2, 1.7) {OS $>$ EN};
    }{Leaving group ability within the periodic table}{fig:filler4}{0.9}

    Comparing down the periodic table for halogens:
    \[
      \ce{F-} < \ce{Cl-} < \ce{Br-} < \ce{I-}
    \]

    \ce{I-} is the best leaving group due to its large orbital size and correspondingly weak bond. Notably, \ce{I-} is such a great leaving group that it does not even participate in radical halogenation, making it impractical in that context. \ce{F-}, on the other hand, has the smallest orbital size, forming much stronger bonds through lower energy orbitals, making it the poorest leaving group of the halogens.

    Comparing across the periodic table (\textit{row 2}):
    \[
      \ce{-CH3} < \ce{-NH2} < \ce{-OH} < \ce{-F}
    \]

    Here, \ce{F-} is the best leaving group due to its greater electronegativity relative to O, N, and C, allowing it to better stabilise the extra \ce{e^-} upon departure. In practice however, \ce{F-} remains a poor leaving group relative to the broader periodic table, \ce{HO-} is exceedingly poor, and \ce{NH2-} and \ce{CH3-} do not leave at all.

    For \textbf{polyatomic} leaving groups, \textbf{resonance} (π conjugation) plays an important additional role. The ability to delocalise \ce{e^-}s across an adjacent p orbital system allows the leaving group to absorb excess charge more effectively --- a larger delocalization system not only stabilises the extra charge but also weakens the leaving group's basicity, making it a better leaving group overall.

    \reactionfig{
      \scheme{
        \chemname{
          \chemfig{
            \textcolor{Green}{\ce{CH3O}}\textcolor{Red}{\ce{^-}}
          }
        }{\textcolor{Blue}{15.5} \quad \textcolor{Green}{\textit{poor}}}
        \qquad
        \chemfig{
          \text{\large<}
        }
        \qquad
        \chemname{
          \chemfig{
            [:0,0.8,,,Green]\textcolor{Green}{\ce{CH3}}
            -[:0]\textcolor{Green}{C}
            (=[:90]\textcolor{Green}{O})
            -[:0]\textcolor{Green}{O}\textcolor{Red}{\ce{^-}}
          }
        }{\textcolor{Blue}{4.7} \quad \textcolor{Green}{\textit{reasonable}}}
        \qquad
        \chemfig{
          \text{\large<}
        }
        \qquad
        \chemname{
          \chemfig{
            [:0,0.8,,,Green]\textcolor{Green}{\ce{CH3}}
            -[:0]\textcolor{Green}{S}
            (=[:90]\textcolor{Green}{O})
            (=[:-90]\textcolor{Green}{O})
            -[:0]\textcolor{Green}{O}\textcolor{Red}{\ce{^-}}
          }
        }{\textcolor{Blue}{-1.2} \quad \textcolor{Green}{\textit{excellent}}}
      }
    }{Polyatomic leaving groups; \textit{corresponding numbers are $pK_a$ values}}{fig:filler5}{1.0}

  \end{core}

  \begin{core}{Nucleophile Quality}

    Nucleophilicity depends on four factors: \textbf{charge-basicity}, \textbf{basicity}, \textbf{solvation}, and \textbf{polarizability}.

    \begin{enumerate}

      \item \textbf{Charge-basicity} refers to the greater tendency of an anion, relative to its neutral counterpart, to form a bond with a cation in order to delocalise its net negative charge.
      \[
        \ce{OH-} > \ce{H2O}\qquad\qquad\ce{NH2-} > \ce{NH3}\qquad\qquad\ce{SO4^{2-}} > \ce{ROSO3-}
      \]

      \item \textbf{Basicity} refers to how readily the nucleophile donates \ce{e^-} density to the electrophile --- analogous to how it would donate to \ce{H+}. As electronegativity decreases moving right to left across the periodic table, nucleophilicity increases accordingly.
      \[
        \ce{H3N} > \ce{H2O}\qquad\qquad\ce{H2N-} > \ce{HO-}\qquad\qquad\ce{HO-} > \ce{F-}
      \]

      \item \textbf{Solvation} refers to the ability of the surrounding solvent to impede the nucleophile's ability to attack the electrophile. In \textbf{polar protic solvents} (\textit{such as water or ethanol, which contain O-H and N-H bonds}), strong hydrogen bonds effectively "handcuff" the nucleophile. Larger nucleophiles are less susceptible to this because their diffuse \ce{e^-} clouds result in lower charge densities, weakening the solvent's grip.
      \[
        \ce{F-} < \ce{Cl-} < \ce{Br-} < \ce{I-}
      \]

      In \textbf{polar aprotic solvents} (\textit{polar solvents lacking O-H and N-H bonds}), hydrogen bonding is absent and the nucleophile is largely unimpeded. The trend therefore reverses to reflect intrinsic basicity.
      \[
        \ce{I-} < \ce{Br-} < \ce{Cl-} < \ce{F-}
      \]

      \textit{Note that \ce{F-} remains a weak base in the grand scheme --- it is simply less solvated than \ce{I-} in protic solvents, and does not compare to stronger bases such as \ce{OH-} or \ce{NH2-}.}

      \item \textbf{Polarizability} refers to the large, diffuse \ce{e^-} clouds of heavier elements, which allow for more effective orbital overlap in the \snt transition state.
      \[
        \ce{H2O} < \ce{H2S} < \ce{H2Se} \qquad\text{or}\qquad \ce{(CH3)3N} < \ce{(CH3)3P}
      \]
    
    \end{enumerate}

    Nucleophilicity arises from a complex interplay of these factors. At its core, it is driven by \textbf{basicity}, which governs the horizontal periodic trend --- atoms become less electronegative and more willing to donate \ce{e^-} density moving left. The vertical trend, however, is governed by the solvent. In \textbf{polar protic solvents}, smaller nucleophiles such as \ce{F-} are heavily solvated, making larger, less-impeded ions like \ce{I-} the stronger attackers. In \textbf{polar aprotic solvents}, these hydrogen-bonding impediments are absent, allowing the intrinsically more basic smaller ions to reassert themselves as the stronger nucleophiles. \textbf{Polarizability} provides larger elements an additional advantage in \snt transition states, as their more diffuse \ce{e^-} clouds can begin overlapping with the electrophilic carbon's orbital from a greater distance.

    \pgfig{

      \miniptable
      \node[font=\small] at (2.5, 1.4) {Periodic Table};

      \draw[->, thick] (4.8, 3.5) -- (0.8, 3.5);
      \node[font=\small] at (2.8, 3.8) {More Basic};

      \draw[->, thick] (5.5, 3.2) -- (5.5, 0.2);
      \node[font=\small, rotate=-90] at (5.8, 1.7) {protic/polarizability};

      \draw[->, thick] (6.4, 0.2) -- (6.4, 3.2);
      \node[font=\small, rotate=-90] at (6.7, 1.7) {aprotic};

    }{Nucleophilic ability within the periodic table}{fig:filler6}{0.9}

  \end{core}

  \begin{core}{Steric Effects on Nucleophilic and Leaving Group Ability}
    
    \textbf{Steric effects} must be considered for both the nucleophile and leaving group. Steric effects arise from \ce{e^-} repulsion between the nucleophile or leaving group and the three substituents bonded to the target electrophile.

    \reactionfig{
      \scheme{
        \chemfig{
          [:0]@{ho1}\ce{HO-}
        }
        \quad
        \chemname{
          \chemfig{
            [:0]\textcolor{Blue}{\ce{CH3}}
            -[:-60]@{carbon1}C
            (<[:-100]H)
            (<:[:-150]H)
            -[:0]\textcolor{Green}{Br}
          }
        }{\textit{PRIMARY attack is easy}}
        \qquad
        \chemfig{
          [:0]@{ho2}\ce{HO-}
        }
        \quad
        \chemname{
          \chemfig{
            [:0]\textcolor{Blue}{\ce{CH3}}
            -[:-60]@{carbon2}C
            (<[:-100]\textcolor{Blue}{\ce{CH3}})
            (<:[:-150]H)
            -[:0]\textcolor{Green}{Br}
          }
        }{\textit{SECONDARY attack is possible}}
        \qquad
        \chemfig{
          [:0]@{ho3}\ce{HO-}
        }
        \quad
        \chemname{
          \chemfig{
            [:0]\textcolor{Blue}{\ce{CH3}}
            -[:-60]@{carbon3}C
            (<[:-100]\textcolor{Blue}{\ce{CH3}})
            (<:[:-150]\textcolor{Blue}{\ce{CH3}})
            -[:0]\textcolor{Green}{Br}
          }
        }{\textcolor{Red}{\textit{TERTIARY attack is impossible}}}
      }
      \chemmovearrow{
        \reactionfigarrow{ho1}{0}{carbon1}{180}{Red}
        \reactionfigarrow{ho2}{0}{carbon2}{180}{Red}
      }
    }{Steric effects in the context of \snt reactions}{fig:filler7}{0.9}

    As a consequence, \textbf{leaving groups are better when larger}, as the repulsion assists their departure, while \textbf{nucleophiles are better when smaller}, as bulk impedes their approach. This steric hindrance is greatest when substituents are directly on the electrophilic carbon (α) and diminishes with distance (β, γ, δ, etc.).

    \reactionfig{
      \scheme{
        \chemfig{
          [:0]
          -[:-30]C\textcolor{Red}{\ce{^δ}}
          (<[:-75])
          (<:[:-105])
          -[:30]C\textcolor{Red}{\ce{^γ}}
          (<[:75])
          (<:[:105])
          -[:-30]C\textcolor{Red}{\ce{^β}}
          (<[:-75])
          (<:[:-105])
          -[:30]\textcolor{Blue}{C}\textcolor{Red}{\ce{^α}}
          (<[:75])
          (<:[:105])
          -[:-30,,,,Green]\textcolor{Green}{L}
        }
      }
    }{Carbon branching}{fig:filler8}{0.9}

    \begin{table}[H]
    \centering
    \begin{tabular}{ccccc}
    \hline
    \textbf{Branch} & \multicolumn{4}{c}{\textbf{Structure \& Relative Rate}} \\
    \hline
     & Hydrogens & Primary & Secondary & Tertiary \\
    α & 145 & 1 & 0.028 & 0 \\
    β & 1 & 0.8 & 0.05 & $10^{-5}$ \\
    \hline
    \end{tabular}
    \caption{Rate of \snt reactions across different branching patterns}
    \label{tab:filler9}
    \end{table} 

  \end{core}

  \begin{core}{\sno}

    \sno reactions are \textbf{unimolecular}, meaning the reaction rate is solely determined by [\ce{R-L}].
    \[
      \text{reaction rate} = k[\ce{R-L}]
    \]

    Unlike \snt, \sno proceeds through a series of distinct steps.
    
    \reactionfig{
      \scheme{
        \chemfig{
          [:0]X
          -[:-60]C
          (<[:-100]Z)
          (<:[:-150]Y)
          -[@{bond}:0]@{leaving}L
        }
        \arrow{->[slow][]}
        \chemfig{
          @{nucleophile}\ce{Nu-}
        }
        \+
        \chembrackets{
          \chemfig{
            [:0]X
            -[:-90]@{carbon}\pcharge{45}{C}
            (<[:-60]Z)
            (<:[:-120]Y)
          }
        }
        \arrow{->[fast][]}
        \chemfig{
          [:0]X
          -[:-60]C
          (<[:-100]Z)
          (<:[:-150]Y)
          -[:0,,,,Blue]\textcolor{Blue}{Nu}
        }
        \+
        \chemfig{
          [:0]\textcolor{Red}{Nu}
          -[:0,,,,Red]C
          (<[:-80]Z)
          (<:[:-30]Y)
          -[:60]X
        }
      }
      \chemmovearrow{
        \reactionfigarrow{nucleophile}{-45}{carbon}{180}{Red}
        \reactionfigarrow{nucleophile}{45}{carbon}{0}{Blue}
      }
    }{\sno mechanism}{fig:filler10}{0.9}

    The first step --- spontaneous dissociation of the leaving group from the electrophile --- is the \textbf{rate limiting step}, as it is highly unfavourable, forming only trace amounts of carbocation at equilibrium. This makes it the bottleneck of the reaction, after which the nucleophile attacks either face of the $sp^2$ hybridised electrophile's perpendicular p orbitals.

    Several important consequences follow from this mechanism.

    \begin{enumerate}

      \item \sno nucleophilic attack is \textbf{not stereospecific}. A stereochemically pure reactant pool will yield a mixture of stereoisomers upon \sno. In practice, \textbf{slight inversion} is observed due to the transient physical blockade created by the lingering leaving group, which remains in proximity to the carbocation p orbital through ionic attraction.

      \item Because the rate is limited solely by carbocation formation, the \textbf{rate increases with a better leaving group}.
      \[
        \ce{-Cl} < \ce{-Br} < \ce{-I} < \ce{-O3SR}
      \]

      \item \textbf{\sno is accelerated by increasingly polar solvents}, particularly polar protic solvents. As the leaving group begins to depart, polar solvent molecules stabilise it through hydrogen bonding, effectively lowering the energy barrier for its departure.

      \doublereactionfig{
        \scheme{
          \chemfig{
            \ce{(CH3)3CBr}
          }
          \quad
          \arrow{->[\textcolor{Red}{100\% \ce{H2O}}][]}
          \quad
          \chemfig{
            \ce{(CH3)3COH}
          }
          \+
          \chemfig{
            \ce{HBr}
          }
          \qquad
          \chemfig{
            \textcolor{Red}{400,000}
          }
        }
      }{
        \scheme{
          \chemfig{
            \ce{(CH3)3CBr}
          }
          \quad
          \qquad
          \arrow{->[\textcolor{Blue}{90\% acetone}][\textcolor{Blue}{10\% \ce{H2O}}]}
          \quad
          \chemfig{
            \ce{(CH3)3COH}
          }
          \+
          \chemfig{
            \ce{HBr}
          }
          \qquad
          \chemfig{
            \textcolor{Blue}{1}
          }
        }
      }{Relative \sno reaction rates in polar solvents}{fig:filler11}{0.9}      

      \item \textbf{Product-determining steps occur after the bottleneck}. Once the carbocation intermediate forms, the dominant product is determined by nucleophilic competition --- the strongest nucleophile in solution will attack preferentially. Consider, for instance, a mixture of NaCl and NaI in acetone.

      \myfig{0.7}{nuc_sub_and_elim_figures/sn1_nu_competition}{sn1 nu competition}{fig:sn1nucomp}
    
    \end{enumerate}

    Here, the overall reaction rate is identical across all paths, as it depends solely on the departure of the leaving group (\textit{regardless of the nucleophile identity}). Because acetone is a polar aprotic solvent, the \ce{Cl-} product will be favored under a mixture of \ce{Cl-} and \ce{I-} because \ce{Cl-} is a stronger nucleophile than \ce{I-} is.

  \end{core}

  \begin{core}{Competition between \snt and \sno}
    
    This raises an important question --- if \sno and \snt ultimately achieve the same outcome, what drives a system to favour one over the other? Two factors determine this.

    \begin{enumerate}


\item α-branching disfavours \snt by sterically impeding nucleophilic approach, making it virtually impossible for a nucleophile to attack a tertiary electrophile in an \snt fashion.

\item α-branched carbocations are stabilised by hyperconjugation with neighbouring molecular orbitals, lowering the energy of the rate-limiting step and making \sno viable.
    
\end{enumerate}

    \myfig{0.75}{nuc_sub_and_elim_figures/sn1_sn2}{sn1 sn2}{fig:sn1sn2}

    The table below summarises the conditions favouring \sno and \snt across different substrate classes. \textit{Note that secondary R substrates do not inherently favour either mechanism --- the outcome depends heavily on solvent and the nature of the nucleophile and leaving group.}

    \begin{center}
      \begin{tabularx}{\textwidth}{lXX}
        \hline
        R & \sno & \snt \\
        \hline
        \ce{CH3} & \textbf{Not observed} (methyl cation too high energy) & \textbf{Frequent} (fast with good Nu/L) \\
        Primary & \textbf{Not observed} (primary carbocation too high energy) & \textbf{Frequent} (fast with good Nu/L, slow when $\beta$ branching present) \\
        Secondary & \textbf{Relatively slow} (best with good L in polar protic solvents) & \textbf{Relatively slow} (best with high concentrations of good Nu in polar aprotic solvents) \\
        Tertiary & \textbf{Frequent} (particularly fast in polar, protic solvents and good L) & \textbf{Extremely slow} \\
        \hline
      \end{tabularx}
    \end{center}

  \end{core}


\section{Elimination}
  
  \begin{core}{Elimination} 

    While substitution reactions involve replacing one group with another, \textbf{elimination} reactions involve \textbf{removal} --- specifically, the loss of a leaving group and an adjacent hydrogen to form a new π bond, yielding an alkene. Like nucleophilic substitution, elimination reactions are classified into two types: E1 and E2.

  \end{core}

  \begin{core}{Elimination 1 (E1)}

    Following carbocation formation in \sno, the nucleophile can alternatively act as a base (\ce{B-}), deprotonating a neighbouring alkyl group to yield a π bonded alkene over two distinct steps. \textit{Note that this pathway requires a neighbouring alkyl group and is therefore not possible for methane.}

    \myfig{0.8}{nuc_sub_and_elim_figures/e1}{e1}{fig:e1}

    The rate-determining step is identical to that of \sno, making E1 a first-order reaction whose rate depends solely on \ce{[R-L]}. A closer look at the deprotonation step is provided below.

    \myfig{}{nuc_sub_and_elim_figures/e1_deprotonation}{e1 deprotonation}{fig:e1depro}

    Since deprotonation can occur at any neighbouring β C-H, E1 generally yields a mixture of alkenes across all possible pathways. Rotation around the resulting π bond gives rise to different stereoisomers.

    \myfig{0.8}{nuc_sub_and_elim_figures/e1_mixture}{e1 mixture}{fig:e1mix}

    Crucially, \sno and E1 share the exact same rate-determining step --- spontaneous departure of the leaving group to form a carbocation. Since the leaving group has already departed by the time the carbocation partitions between substitution and elimination, the \textbf{\sno to E1 product ratio is independent of L}.

    The ratio between E1 and \sno products is difficult to predict precisely, but E1 is generally favoured under two conditions.

    \begin{enumerate}


\item \textbf{Higher temperature}, as elimination carries a positive entropy change: \ce{HRL + B -> R + L + BH}.

\item A \textbf{poor nucleophile}, as \sno requires the reagent to form a new C-Nu bond, demanding greater nucleophilic strength. E1 is less sensitive to this since the reagent acts purely as a base, abstracting a surface \ce{H+} rather than integrating into the molecular framework.

    
\end{enumerate}

    But what about \textbf{strong bases}?

  \end{core}

  \begin{core}{Elimination 2 (E2)}

    Strong bases attack \ce{R-L} \textbf{directly} at the β C-H, constituting the E2 mechanism. E2 is faster than \sno/E1 and faster than \snt when sterically hindered by α or β branching, as it condenses the distinct steps of E1 into a single concerted process.

    \myfig{0.9}{nuc_sub_and_elim_figures/e2_ts}{e2 ts}{fig:e2ts}

    Unlike E1, E2 is \textbf{second order} because the transition state involves both the substrate and the base.
    \[
      Rate = k[\ce{R-L}][\ce{B-}]
    \]

    Consequently, a better leaving group and a stronger base both accelerate the reaction.
    \[
      \ce{RCl} < \ce{RBr} < \ce{RI}
    \]

    Additionally, electron-withdrawing substituents within the substrate accelerate E2 by increasing the acidity of the target β C-H through withdrawal of its \ce{e^-} density.

    \myfig{0.5}{nuc_sub_and_elim_figures/e2_en}{e2 en}{fig:e2en}

    E2 is also \textbf{stereospecific} --- a given diastereomer of \ce{R-L} yields only one stereoisomer of the alkene product. For instance, an RR/SS mixture yields exclusively one product, while an RS/SR mixture yields exclusively another.

    \myfig{0.6}{nuc_sub_and_elim_figures/rrss}{rrss}{fig:rrss}

    \myfig{0.6}{nuc_sub_and_elim_figures/rssr}{rssr}{fig:rssr}

    This stereospecificity arises because \ce{H+} and L depart simultaneously, requiring their initial positions to be coplanar to permit proper p orbital alignment in the forming $\pi$ bond. This coplanarity can be achieved through either \textbf{anti}-periplanar (\textit{trans}) or \textbf{syn}-periplanar (\textit{cis}) geometry. When the relevant C-C bond can rotate freely, anti-periplanar geometry is strongly preferred as it minimises steric repulsion between the leaving group and the \ce{H-B} transition complex.

    In cases where free rotation is restricted --- such as in cycloalkanes --- syn-periplanar geometry must be adopted instead. This proceeds at a slower rate due to the increased steric hindrance between the leaving group and the \ce{H-B} transition complex. 

    \myfig{0.65}{nuc_sub_and_elim_figures/anti_syn}{anti syn}{fig:antisyn}

  \end{core}
  
\section{Competition}

  Having established the mechanics and characteristics of \snt, \sno, E1, and E2 reactions, it is important to recognise that all four pathways can be simultaneously possible across many overlapping conditions --- and thus compete with one another in a given reaction environment. The table below organises these competing factors.

  \begin{table}[H]
  \centering
  \footnotesize
  \begin{tabularx}{\textwidth}{lXXXX}
  \hline
  \textbf{Type of haloalkane} & \textbf{Poor Nu (e.g. \ce{H2O})} & \textbf{Weakly basic, good Nu} & \textbf{Strongly basic, unhindered Nu} & \textbf{Strongly basic, hindered Nu} \\
  \hline
  Methyl & N/A & S$_N$2 & S$_N$2 & S$_N$2 \\
  Primary & & & & \\
  \quad Unhindered & N/A & S$_N$2 & S$_N$2 & E2 \\
  \quad Branched & N/A & S$_N$2 & E2 & E2 \\
  Secondary & Slow S$_N$1, E1 & S$_N$2 & E2 & E2 \\
  Tertiary & S$_N$1, E1 & S$_N$1 & E2 & E2 \\
  \end{tabularx}
  \end{table}

  \begin{core}{Factors of Competition between \snt, \sno, E1, and E2}

    Three core factors govern the competition between substitution and elimination.

    \begin{enumerate}


\item The base strength of the nucleophile

\item Steric hindrance around the reacting carbon

      \item Steric hindrance around the nucleophile, particularly when it is a strong base

    
\end{enumerate}

    \textbf{Primary substrates:} Under poor nucleophile conditions, neither E1 nor \sno are viable due to insufficient hyperconjugation to sustain a carbocation, and neither \snt nor E2 proceed effectively as both depend on nucleophile strength. Moving from weakly basic to strongly hindered nucleophile conditions, the dominant factor becomes steric hindrance --- whether along the haloalkane itself via β branching, or around the nucleophile through size or solvation in protic solvents. As hindrance increases, E2 is progressively favoured over \snt.

    \textbf{Secondary substrates:} Hyperconjugation is sufficient to support carbocation formation, making \sno and E1 viable under poor nucleophile conditions. E1 tends to dominate over \sno here, as abstracting a β \ce{H+} is less demanding than integrating a nucleophile into the molecular framework. For stronger nucleophile conditions, the inherent steric hindrance of the secondary substrate generally drives the system toward E2, though \snt remains competitive when the nucleophile is weakly basic --- insufficiently basic to abstract the β \ce{H+} required by the E2 transition state.

    \textbf{Tertiary substrates:} Steric hindrance completely shuts down \snt, making \sno dominant under weakly basic, good nucleophile conditions. E2 dominates under strongly basic conditions, as it does for secondary substrates. Under poor nucleophile conditions, the \sno to E1 ratio becomes more balanced relative to secondary substrates --- the greater stability of the tertiary carbocation lowers the activation energy for \sno, allowing it to compete more effectively with E1.

    Generally speaking, rather than memorizing each of these scenarios and their corresponding dominant mechanism, it's helpful to start with \snt as a starting point and ask yourself the following questions.

    \begin{enumerate}


\item Is the nucleophile or haloalkane hindered? \textit{If yes, the \snt reaction is likely not favored.}

\item Is the nucleophile strongly basic? \textit{If yes, E2. If no, \sno/E1.}

    
\end{enumerate}

    However, keep in mind that there are other nuanced factors to consider such as absence of E1 preference for weakly basic, good nucleophile conditions.

  \end{core}
  
\end{document}

